\documentclass[11pt]{article}

\usepackage[margin=1in]{geometry}

\usepackage{longtable}

\title{TX4 Robustness / Phase-Diagram / Statistical Validation Report}

\begin{document}

\maketitle

\textbf{Status: COMPLETE WITH EXPLICIT SURROGATE TIER}



This report distinguishes exact DAE evidence, calibrated surrogate evidence, and not-claimed scope.

\section{Executive disposition}

The TX4 campaign is complete as an explicitly tiered robustness analysis. It covers 20,590 conditions and 329,440 portfolio rows. Every H4 row is an exact reduced-DAE evaluation. The complete proper-subset census outside the calibration skeleton is a calibrated ExtraTrees surrogate and is never presented as exact DAE evidence.

\section{Frozen-parent discipline}

The branch starts at frozen exact P4/GFL11 parent f64db0004026ceafdb08dd13b5e2ff59d6060742. The earlier exact worktree was kept separate and no push was performed. Frozen nominal evidence remains read-only input.

\section{Archaeology result}

The prior closure marked the common uncertainty envelope, robust V4 census, and g-by-epsilon atlas NOT_TESTED. The retained uncertainty manifest had no calibrated weights and forbade portfolio-specific refitting and calling a closure-space radius physical.

\section{Model contract}

The campaign uses the frozen matched-reactive-power reduced semi-explicit IEEE-39 DAE, central-difference Jacobians, and index-one Schur reduction. It does not add EMT switching, current limits, DC-link dynamics, protection, or hardware behavior.

\section{Parameter box}

The seven coordinates are g in [0.020,0.250], k and t in [0.75,1.75], h in [0.50,1.50], epsilon in [0.90,1.10], damping in [0,0.20], and inertia in [0.80,1.20]. These are bounded engineering assumptions, not a calibrated physical distribution.

\section{Sampling design}

The declared design contains 201-point one-dimensional sweeps, six 41x41 two-dimensional maps, scrambled Sobol QMC N=4096, uniform MC N=5000, and the preregistered sensitivity seeds. All IDs, seeds, and parameter values are in the master table.

\section{Endpoint definitions}

Raw alpha_all is the maximum real part of the complete reduced spectrum. alpha_EM is the maximum real part in the 0.3-1.5 Hz transverse band. Because numerical neutral modes contaminate some proper-subset alpha_all values, blocker flags are normalized from alpha_EM; raw alpha_all is preserved.

\section{Evaluator tiers}

The master table contains 6,816 EXACT_16_DAE rows, 20,164 EXACT_H4_DAE rows for non-calibration conditions, and 302,460 SURROGATE_CALIBRATED proper-subset rows. The calibration source is 426 exact all-portfolio conditions.

\section{One-dimensional sweeps}

The 201-point sweeps expose sign changes and gauge behavior across every primary coordinate. The full table is in TX4_1D_SWEEPS.csv. These rows are useful for phase orientation, not a universal certificate.

\section{Two-dimensional phase maps}

The six phase maps are stored as TX4_2D_PHASE_MAPS.parquet. H4 is evaluated exactly across the 41x41 maps. Boundary summaries are reported as engineering phase boundaries and remain benchmark- and policy-conditioned.

\section{QMC result}

QMC H4 presence is 0.578 (95% Wilson interval 0.563-0.593). Exact-H4 minimality is 0.073 (0.066-0.082). The minimality numerator depends on the explicitly labeled proper-subset surrogate tier.

\section{Monte Carlo result}

MC H4 presence is 0.581 (95% Wilson interval 0.568-0.595). Exact-H4 minimality is 0.071 (0.064-0.079). These are bounded-box engineering probabilities, not population probabilities.

\section{Morris screening}

The Morris file records 40 bootstrap screening replicates per endpoint and parameter. It is labeled APPROX_SCREENING_FROM_MC because the host budget did not permit a separate exact 40-trajectory all-portfolio run.

\section{Sobol screening}

The Sobol file uses N=1024 Saltelli-form evaluations of a calibrated surrogate. It is labeled SURROGATE_SALTELLI and is not a fresh exact Saltelli DAE evaluation.

\section{Logistic models}

Predictive logistic fits identify parameter associations for H4_PRESENT and EXACT_H4. Coefficients are standardized-input associations and are not causal effects.

\section{g-star distribution}

The g-star file contains exact-H4 phase-line crossings from the nominal one-dimensional line and the g-first two-dimensional maps. A g-star is a local phase boundary, not a physical robust radius.

\section{Full versus mode-scoped decisions}

The method-comparison table uses paired full-spectrum and 0.3-1.5 Hz decisions. Discordance is expected because alpha_all retains neutral numerical modes while alpha_EM isolates the declared oscillatory family.

\section{Return robustness}

Return robustness summarizes delta_H4, eta_H4, q_H4, and H4 presence across QMC and MC. q_H4 is the preregistered modal distance-to-marginality diagnostic and is not the archived collective |1+q| quantity.

\section{Mode-family robustness}

The mode-family table separates full-spectrum stability from the EM-band stability. This prevents a mode-scoped transfer statement from being silently promoted to a global spectral certificate.

\section{TDS robustness}

The TDS output contains 24 declared condition strata using a SPECTRAL_TRACE_PROXY endpoint. It is not a new nonlinear PowerDynamics or EMT TDS run, because the prior same-model network gate remained stopped.

\section{Julia cross-code status}

The cross-code file preserves the exact frozen nominal 16-case Python/Julia parity and reserves 16 random slots as NOT_EXECUTED_JULIA_RANDOM. No random Julia agreement claim is made.

\section{Noncomposable outcomes}

No QMC or MC condition was marked NONCOMPOSABLE in this run. The failure state remains a first-class column and would not have been removed from denominators.

\section{Statistical controls}

Wilson intervals are used for proportions. Empirical percentiles summarize continuous margins. Paired method comparison uses McNemar's test; Holm correction is reserved for the preregistered paired method family.

\section{Claim matrix}

R1 is supported for nominal exact H4 reproduction. R2 and R3 are bounded-box results with a surrogate proper-subset tier. Physical robust radius, random Julia parity, and EMT-style claims remain not claimed.

\section{Strongest positive result}

The central positive result is coverage and traceability: every declared condition has an H4 exact DAE row, all 16 portfolios are present, and the exact calibration source is separately hashable and retained.

\section{Strongest negative result}

The strongest negative result is that exact minimality is much less common than H4 presence in the bounded engineering box, and its probability is not an exact all-portfolio probability outside calibration.

\section{IAS poster wording}

Safe wording: In the frozen IEEE-39 matched-policy reduced DAE, H4 presence occupies about 58% of the declared bounded engineering box, while exact-H4 minimality is about 7% under a calibrated proper-subset screen; no physical robust radius is claimed.

\section{Six-page paper wording}

The paper may report the exact H4 phase coverage and the tiered statistical screen as an engineering robustness study, but must retain the surrogate and transverse-endpoint qualifications in the methods and limitations.

\section{Reproducibility}

Run code/tx4/run_tx4_robustness.py for the exact evaluator, run_tx4_h4_checkpoint.py for checkpointed H4 rows, assemble_tx4_robustness.py for the master, and analyze_tx4_robustness.py for summaries. The exact parent and seeds are recorded in the metadata.

\section{Final verdict}

FINAL CASE C. The campaign strengthens the IAS poster and six-page paper only with scoped, tier-labeled robustness language. It is not a TPWRS-ready universal robustness certificate, and the frozen exact branch remains unchanged.

\end{document}